\begin{beginnerstatement}[Kurz erklärt: Schwäche, Grenze und Nachweis]

Eine technische Schwäche ist ein tatsächlicher Nachteil der vorhandenen
Umsetzung.

Eine Scope-Grenze legt dagegen bewusst fest, was das Template nicht abdecken
soll.

Fehlende Verifikation bedeutet, dass eine Funktion noch nicht ausreichend
geprüft oder nachgewiesen wurde; sie ist deshalb nicht automatisch fehlerhaft.

Eine Regression ist ein Fehler, der durch eine Änderung in einem zuvor
funktionierenden Bereich entsteht.

Eine Regel kann außerdem technisch umgangen werden, obwohl die Dokumentation
dies nicht vorsieht. Im untersuchten Stand kann eine Exception sogar den
901-LOC-Fehler unterdrücken. Das ist eine technische Schwäche und keine
Bestätigung, dass 901 LOC erlaubt sein sollen. Die Fallstudie nennt daher
sowohl die beabsichtigte Regel als auch das tatsächlich beobachtete Verhalten.

\end{beginnerstatement}
